\subsection{What is inherited}

JHQ combines an orthogonal transformation, an analytical Cartesian primary
code and a higher-precision residual layer \cite{jhq}. Here the transform
is square: $\Pi\in\mathbb R^{d\times d}$, $\Pi^\top\Pi=I$,
$y=\Pi x$ and $q'=\Pi q$. Therefore
\begin{equation}
 \|q-x\|^2=\|\Pi(q-x)\|^2.
 \label{eq:2-isometry}
\end{equation}
Although the prior work calls this an orthogonal JL transformation, this
configuration does not reduce dimension. It preserves real-arithmetic
distance exactly; quantisation subsequently introduces approximation.
Actual FP32 orthogonality and GEMM error are separate implementation checks.

Partition $y$ into $M$ subspaces of dimension $D_s=d/M$. With $B$ bits per
primary subspace, $K=2^B$ labels are formed from $L=2^{B/D_s}$ scalar levels,
requiring the implemented Cartesian parameters to make $L$ an integer.
Writing $\Phi$ for the standard normal CDF, the constructor uses
\begin{equation}
 a_i=\sigma\Phi^{-1}\!\left(\frac{i-1/2}{L}\right),\quad
 i=1,\ldots,L,\qquad
 \mathcal C_m=\{a_1,\ldots,a_L\}^{D_s}.
 \label{eq:2-cartesian}
\end{equation}
At the evaluated $B=D_s=8$, $L=2$ and $\mathcal C_m=\{-a,+a\}^8$,
with $a\simeq0.67449\sigma$. Each byte records eight coordinate signs.
The implementation estimates $\sigma$ from vector energy; unit-normalised
inputs motivate the scale $1/\sqrt d$. This representation is prior art,
not an invention of the GPU execution layer.

Residuals are $r=y-\widehat p$, giving reconstruction $\widehat y=\widehat p+\widehat r$. The evaluated codes use $d/8$ primary bytes and $d$ residual bytes per vector, before metadata. Refining fewer rows reduces reads, not stored residual size.

\subsection{Model assumptions and score semantics}

Random rotation can spread coordinate energy, motivating a shared
Gaussian-inspired scalar model. It does not force an arbitrary realised
database to have independent, zero-mean Gaussian coordinates. For a fixed
rotation and empirical database mean $\mu$, the transformed mean is
$\Pi\mu$ and $\|\Pi\mu\|=\|\mu\|$. Averaging over random rotations is
different from centring one fixed corpus. Even a uniformly random point on
a sphere has dependent coordinates through its fixed norm. Consequently,
isometry is an algebraic guarantee; the distributional approximation is a
modelling assumption. Chapter 8's centred-coding pilot examines one practical
consequence of that distinction.

The inherited primary levels are Gaussian mid-quantiles. Under an exact Gaussian model, zero-threshold one-bit MSE-optimal levels would instead be $\pm\sigma\sqrt{2/\pi}\simeq\pm0.79788\sigma$. For fixed signs and positive scale, however,
\begin{equation}
 \primarydist=\|q'\|^2+da^2-2a\langle q',s\rangle,
 \qquad s\in\{-1,+1\}^d.
 \label{eq:2-sign-order}
\end{equation}
Thus scale cancels from primary ordering but still affects residual reconstruction. Our design uses separability and measured quality, without claiming optimal quantisation.

$\rawdist$ uses original vectors; $\primarydist=\|q'-\widehat p\|^2$ uses primary reconstructions. The refined score is
\begin{equation}
 \refineddist=\|q'-\widehat p-\widehat r\|^2+\|q'\|^2.
 \label{eq:2-refined}
\end{equation}
Its query-only term changes distance values, not ranking. Routing and reconstruction errors remain after exact computed-score selection.